\section*{अध्याय \ref{ch:g12:reaction-rates} --- अभिक्रिया दरें}

\begin{solution}{exo:g12:reaction-rates:1}
ताप; संपर्क क्षेत्रफल; सांद्रता; ताप।
\end{solution}

\begin{solution}{exo:g12:reaction-rates:2}
\qty{13.9}{min} (नीला) और \qty{6.9}{min} (लाल)। दो \omterm{def:g12:reaction-rates:half-life}{अर्ध-आयुओं} के बाद,
$10.0 / 4 = \qty{2.5}{mmol/L}$।
\end{solution}

\begin{solution}{exo:g12:reaction-rates:3}
प्रवणता $(5.0 - 1.0)/10 = \qty{0.40}{mmol.L^{-1}.min^{-1}}$: यही
\qty{5}{min} पर \omterm{def:g12:reaction-rates:rate}{प्रकटन दर} है।
\end{solution}

\begin{solution}{exo:g12:reaction-rates:4}
ठंडा और तनु करने से अभिक्रिया इतनी धीमी हो जाती है कि \omterm{def:g11:titration:titration}{अनुमापन} के दौरान संरचना
अब नहीं बदलती। मापी गई संरचना वही है जो
\omterm{def:g6:pure-substances-and-mixtures:mixture}{मिश्रण} की तब थी जब अभिक्रिया रोकी गई, डालने के क्षण।
\end{solution}

\begin{solution}{exo:g12:reaction-rates:5}
\qty{30}{min} 3 \omterm{def:g12:reaction-rates:half-life}{अर्ध-आयु} है: $1/8$ बचता है। \qty{1}{h} 6 \omterm{def:g12:reaction-rates:half-life}{अर्ध-आयु} है:
$1/64$, लगभग \qty{1.6}{\percent}।
\end{solution}

\begin{solution}{exo:g12:reaction-rates:6}
$\ln[\ce{A}]$: 2.079, 1.733, 1.386, 1.040, 0.693। यह हर \qty{5}{min} पर
0.346 घटता है: एक सीधी रेखा, प्रथम कोटि। $k = 0.346/5 =
\qty{0.0693}{min^{-1}}$ और $t_{1/2} = \ln 2 / k = \qty{10.0}{min}$ (जैसा
सीधे दिखता है: हर \qty{10}{min} पर 8.00, 4.00, 2.00)।
\end{solution}

\begin{solution}{exo:g12:reaction-rates:7}
$k = 0.693 / 25 = \qty{0.0277}{s^{-1}}$।
\end{solution}

\begin{solution}{exo:g12:reaction-rates:8}
$0.100\,e^{-0.20} = \qty{0.0819}{mol/L}$; $0.100\,e^{-0.60} =
\qty{0.0549}{mol/L}$; $0.100\,e^{-2.0} = \qty{0.0135}{mol/L}$।
\end{solution}

\begin{solution}{exo:g12:reaction-rates:9}
\qty{13.9}{min} पर $[\ce{A}] = \qty{5.0}{mmol/L}$; स्पर्शरेखा की प्रवणता
लगभग $-0.25$ है, और $-k[\ce{A}] = -0.050 \times 5.0 =
\qty{-0.25}{mmol.L^{-1}.min^{-1}}$: प्रारंभिक मान का आधा।
\end{solution}

\begin{solution}{exo:g12:reaction-rates:10}
आयोडीन ही एकमात्र रंगीन स्पीशीज़ है: बियर--लैम्बर्ट नियम से
$A = \varepsilon \ell [\ce{I2}]$। \omterm{def:g12:reaction-rates:rate}{प्रकटन दर} वक्र $[\ce{I2}](t)$ की
स्पर्शरेखा की प्रवणता है, अर्थात $A(t)$ की
स्पर्शरेखा की प्रवणता को $\varepsilon \ell$ से भाग देने पर।
\end{solution}

\begin{solution}{exo:g12:reaction-rates:11}
प्रारंभिक दर $k[\ce{A}]_0$ दूसरे प्रयोग के लिए दोगुनी है;
\omterm{def:g12:reaction-rates:half-life}{अर्ध-आयु} बराबर हैं, $\ln 2 / k$, जो $[\ce{A}]_0$ पर निर्भर नहीं।
\end{solution}

\begin{solution}{exo:g12:reaction-rates:12}
$t = \ln(1/f)/k$। \omterm{def:g12:reaction-rates:half-life}{अर्ध-आयुओं} में: $\ln(1/f)/\ln 2$। \qty{1}{\percent} के लिए
$\ln 100 / \ln 2 = 6.6$; \qty{0.1}{\percent} के लिए $\ln 1000 / \ln 2 =
10.0$।
\end{solution}

\begin{solution}{exo:g12:reaction-rates:13}
\qty{90}{\percent} खपत, $f = 0.10$: $t = \ln 10 / k$, अर्थात
\qty{20}{\celsius} पर \qty{46}{min} और \qty{30}{\celsius} पर \qty{23}{min}।
\omterm{def:g12:reaction-rates:first-order}{दर स्थिरांक} \qty{10}{\celsius} पर दोगुना हो जाता है: यहाँ नियम लागू होता है।
\end{solution}

\begin{solution}{exo:g12:reaction-rates:14}
ताप लगभग \qty{0}{\celsius} तक गिर जाता है, और सांद्रताएँ
10 से विभाजित हो जाती हैं। दो सांद्रताओं के गुणनफल के समानुपाती दर के साथ,
केवल \omterm{def:g10:concentration-and-dilution:dilution}{तनुकरण} उसे $10 \times 10 = 100$ से विभाजित करता है।
\end{solution}

\begin{solution}{exo:g12:reaction-rates:15}
$v(t) = -\dfrac{d}{dt}\big([\ce{A}]_0 e^{-kt}\big) = k[\ce{A}]_0 e^{-kt}$।
प्रारंभिक दर $k[\ce{A}]_0$; दर तब आधी होती है जब $e^{-kt} = 1/2$, अर्थात
एक \omterm{def:g12:reaction-rates:half-life}{अर्ध-आयु} के बाद।
\end{solution}

\begin{solution}{pb:g12:reaction-rates:1}
\textbf{1.} बियर--लैम्बर्ट नियम से, क्योंकि रंजक ही इस तरंगदैर्घ्य पर अवशोषित करने वाली
एकमात्र स्पीशीज़ है; विलयन पर्याप्त तनु होना चाहिए
($A$ लगभग 1 से कम)।

\textbf{2.} वह कोई दृश्य प्रकाश अवशोषित नहीं करता: वह $A$ में कुछ नहीं जोड़ता।

\textbf{3.} ताकि उसकी सांद्रता मुश्किल से बदले: तब दर
केवल रंजक पर निर्भर करती है।

\textbf{4.} $A = 0$ (रिक्त विलयन)।

\textbf{5.} $A$ \qty{6}{min} में 0.800 से 0.402 तक गिरता है: लगभग
\qty{6}{min}।

\textbf{6.} \qty{6}{min} पर 0.402 से \qty{12}{min} पर 0.199 तक: फिर से
\qty{6}{min} में आधा।

\textbf{7.} $-0.223$, $-0.448$, $-0.691$, $-0.911$, $-1.155$, $-1.366$,
$-1.614$, $-1.945$, $-2.551$, $-3.079$।

\textbf{8.} बिंदु एक सीधी रेखा पर हैं: अभिक्रिया रंजक के सापेक्ष प्रथम
कोटि की है।

\textbf{9.} $(-2.551 + 0.223)/20 = \qty{-0.116}{min^{-1}}$।

\textbf{10.} $k = \qty{0.116}{min^{-1}}$।

\textbf{11.} $0.693 / 0.116 = \qty{6.0}{min}$, जैसा सारणी में पढ़ा गया।

\textbf{12.} $k A_0 = 0.116 \times 0.800 = \qty{0.093}{min^{-1}}$,
\omterm{def:g11:absorbance:absorbance}{अवशोषकता} की इकाइयों में प्रति मिनट।

\textbf{13.} वही, \qty{6.0}{min}: प्रथम कोटि की \omterm{def:g12:reaction-rates:half-life}{अर्ध-आयु} प्रारंभिक मान
पर निर्भर नहीं करती।

\textbf{14.} $0.800\, e^{-0.116 \times 30} = 0.025$।

\textbf{15.} $t = \ln 100 / k = 4.61 / 0.116 = \qty{40}{min}$।

\textbf{16.} $40 / 6.0 = 6.6$ \omterm{def:g12:reaction-rates:half-life}{अर्ध-आयु}।

\textbf{17.} $t = \ln(0.800/0.010)/k = 4.38 / 0.116 = \qty{38}{min}$।

\textbf{18.} फिर भी 6.6 \omterm{def:g12:reaction-rates:half-life}{अर्ध-आयु}, अब $6.6 \times 3.0 = \qty{20}{min}$।

\textbf{19.} लगभग 6.6 \omterm{def:g12:reaction-rates:half-life}{अर्ध-आयु}, अर्थात यहाँ लगभग \qty{40}{min}।
\end{solution}
